\documentclass{article}
\usepackage[margin=1in, a4paper]{geometry}
\usepackage{amsmath, amssymb, amsfonts}
\usepackage{graphicx}
\usepackage{xcolor}
\usepackage{listings}
\usepackage{booktabs}
\usepackage[utf8]{inputenc}
\usepackage[T1]{fontenc}
\usepackage{hyperref}
\hypersetup{colorlinks=true, urlcolor=blue, linkcolor=blue}
\usepackage{enumitem}
\usepackage{float}

% Professional CMYK color palette
\definecolor{myblue}{cmyk}{0.8,0.4,0,0.1}
\definecolor{mygreen}{cmyk}{0.7,0,0.5,0.2}
\definecolor{myorange}{cmyk}{0,0.5,0.8,0.1}
\definecolor{mygray}{cmyk}{0,0,0,0.4}
\definecolor{lightcyan}{cmyk}{0.3,0,0,0.05}
\definecolor{lightorange}{cmyk}{0,0.3,0.5,0.1}

\definecolor{myblue}{cmyk}{0.8,0.4,0,0.1}
\definecolor{mygreen}{cmyk}{0.7,0,0.5,0.2}
\definecolor{myorange}{cmyk}{0,0.5,0.8,0.1}
\definecolor{mygray}{cmyk}{0,0,0,0.4}
\definecolor{arrowcolor}{cmyk}{0.9,0.5,0,0.3}
\definecolor{nodecolor}{cmyk}{0.6,0.2,0,0.05}
\definecolor{framecolor}{cmyk}{0.4,0.1,0,0.1}

\definecolor{codegreen}{rgb}{0,0.6,0}
\definecolor{codegray}{rgb}{0.5,0.5,0.5}
\definecolor{codepurple}{rgb}{0.58,0,0.82}
\definecolor{backcolour}{rgb}{0.99,0.99,0.99}

\lstdefinestyle{mystyle}{
    backgroundcolor=\color{backcolour},   
    commentstyle=\color{codegreen},
    keywordstyle=\color{magenta},
    numberstyle=\tiny\color{codegray},
    stringstyle=\color{codepurple},
    basicstyle=\ttfamily\footnotesize,
    breakatwhitespace=false,         
    breaklines=true,                 
    captionpos=b,                    
    keepspaces=true,                 
    numbers=left,                    
    numbersep=5pt,                  
    showspaces=false,                
    showstringspaces=false,
    showtabs=false,                  
    tabsize=2
}
\lstset{style=mystyle}

\title{The EOQ with Planned Shortages Problem}
\author{The Logistics \& Supply Chain Problem-Solving Playbook}
\date{}

\begin{document}

\maketitle

\section{The EOQ with Planned Shortages Problem}

\subsection{The Scenario: Strategic Inventory Management with Controlled Stockouts}

In the highly competitive landscape of modern supply chain management, the traditional Economic Order Quantity (EOQ) model's assumption of zero stockouts often proves too restrictive and economically suboptimal. The Meridian Electronics Distribution Center faces this exact challenge while managing semiconductor components for major technology manufacturers.

The distribution center supplies critical microprocessors with an annual demand of 36,000 units, experiencing significant seasonal fluctuations and supply chain uncertainties. While the traditional EOQ model would require maintaining sufficient inventory to never experience stockouts, Meridian has discovered that their customers are increasingly willing to accept planned shortages in exchange for lower overall costs, particularly for non-critical components with predictable delivery windows.

The company faces a complex cost structure: each purchase order costs \$150 to process, annual holding costs amount to \$8 per unit per year, and shortage costs are estimated at \$12 per unit per year when customers must wait for backorders. The key insight is that by strategically allowing controlled shortages during certain periods, the total system costs can be minimized while maintaining acceptable service levels.

This scenario represents the EOQ with Planned Shortages Problem, where the objective shifts from eliminating stockouts entirely to finding the optimal balance between inventory holding costs, ordering costs, and shortage costs. The solution determines not only the optimal order quantity but also the planned shortage level that minimizes total system costs.

\subsection{Tier 1: The Pen \& Paper Method (Mixed-Integer Programming Formulation)}

The EOQ with Planned Shortages problem can be elegantly formulated as a mixed-integer programming model that captures the fundamental trade-offs between inventory costs and shortage costs.

\textbf{Mathematical Model Formulation}

Let us define the following sets, parameters, and decision variables:

\textbf{Parameters:}
\begin{align}
D &= \text{Annual demand rate (units/year)} \\
K &= \text{Fixed ordering cost per order (\$/order)} \\
h &= \text{Holding cost per unit per year (\$/unit/year)} \\
p &= \text{Shortage cost per unit per year (\$/unit/year)} \\
T &= \text{Planning horizon (years)}
\end{align}

\textbf{Decision Variables:}
\begin{align}
Q &= \text{Order quantity (units)} \\
S &= \text{Maximum shortage level (units)} \\
I_{max} &= \text{Maximum inventory level (units)} \\
t_1 &= \text{Time with positive inventory (years)} \\
t_2 &= \text{Time with shortages (years)} \\
n &= \text{Number of orders per planning horizon (integer)}
\end{align}

The cycle relationships establish that $I_{max} = Q - S$ and the cycle time is $t = t_1 + t_2 = \frac{Q}{D}$.

\textbf{Objective Function:}

The total cost per planning horizon consists of three components:

\begin{align}
\text{Minimize } TC &= \text{Ordering Cost} + \text{Holding Cost} + \text{Shortage Cost} \\
TC &= n \cdot K + \frac{I_{max}^2 \cdot t_1 \cdot h}{2Q} + \frac{S^2 \cdot t_2 \cdot p}{2Q}
\end{align}

Substituting the relationships $I_{max} = Q - S$, $t_1 = \frac{Q-S}{D}$, $t_2 = \frac{S}{D}$, and $n = \frac{D \cdot T}{Q}$:

\begin{equation}
TC = \frac{D \cdot T \cdot K}{Q} + \frac{(Q-S)^2 \cdot h}{2D} + \frac{S^2 \cdot p}{2D}
\end{equation}

\textbf{Constraints:}
\begin{align}
Q &\geq S \geq 0 \\
Q &> 0 \\
n &\geq 1, \text{ integer}
\end{align}

\textbf{Optimal Solution Derivation:}

Taking partial derivatives and setting them to zero:

\begin{align}
\frac{\partial TC}{\partial Q} &= -\frac{D \cdot T \cdot K}{Q^2} + \frac{h(Q-S)}{D} = 0 \\
\frac{\partial TC}{\partial S} &= -\frac{h(Q-S)}{D} + \frac{pS}{D} = 0
\end{align}

Solving simultaneously yields the optimal solutions:

\begin{align}
Q^* &= \sqrt{\frac{2DTK(h+p)}{hp}} \\
S^* &= Q^* \cdot \frac{h}{h+p} \\
I_{max}^* &= Q^* \cdot \frac{p}{h+p}
\end{align}

\textbf{Inventory Level Visualization:}

\begin{figure}[htbp]
\centering
\includegraphics[width=\textwidth]{../figures-png/line53.fig1.png}
\caption{Enhanced Quantity Discount EOQ Problem with Professional CMYK Colors and Node-Based Cost Analysis}
\end{figure}

\textbf{Concrete Numerical Example:}

Consider Meridian Electronics with the following parameters:
\begin{align}
D &= 36,000 \text{ units/year} \\
K &= \$150 \text{ per order} \\
h &= \$8 \text{ per unit per year} \\
p &= \$12 \text{ per unit per year}
\end{align}

Calculating the optimal solution:

\begin{align}
Q^* &= \sqrt{\frac{2 \times 36,000 \times 150 \times (8+12)}{8 \times 12}} = \sqrt{\frac{216,000,000}{96}} = 1,500 \text{ units} \\
S^* &= 1,500 \times \frac{8}{8+12} = 1,500 \times 0.4 = 600 \text{ units} \\
I_{max}^* &= 1,500 \times \frac{12}{8+12} = 1,500 \times 0.6 = 900 \text{ units}
\end{align}

The cycle time is $T = \frac{Q^*}{D} = \frac{1,500}{36,000} = 0.0417$ years (approximately 15.2 days).

Time with inventory: $t_1 = \frac{I_{max}^*}{D} = \frac{900}{36,000} = 0.025$ years (9.1 days)

Time with shortages: $t_2 = \frac{S^*}{D} = \frac{600}{36,000} = 0.0167$ years (6.1 days)

Total annual cost: $TC = \frac{36,000 \times 150}{1,500} + \frac{900^2 \times 8}{2 \times 36,000} + \frac{600^2 \times 12}{2 \times 36,000} = \$3,600 + \$90 + \$60 = \$3,750$

\subsection{Tier 2: The Classic Heuristic (Greedy Construction Algorithm)}

While the mathematical formulation provides the theoretical optimum, practical implementations often require heuristic approaches that can handle real-world complexities such as discrete order quantities, supplier constraints, and dynamic demand patterns.

\textbf{Algorithm Design Philosophy}

The greedy heuristic approach constructs a feasible solution by making locally optimal decisions at each step. The algorithm incrementally adjusts the order quantity and shortage level based on marginal cost improvements, ensuring that each modification reduces the total system cost.

\textbf{Greedy EOQ with Shortages Algorithm}

\lstinputlisting[language=Python]{.././codes-python/line53.py1.py}

\textbf{Algorithm Visualization}

\begin{figure}[htbp]
\centering
\includegraphics[width=\textwidth]{../figures-png/line53.fig2.png}
\caption{Greedy Algorithm Flow for EOQ with Shortages}
\end{figure}

\textbf{Concrete Example Results}

Running the greedy algorithm on the Meridian Electronics example:

\lstinputlisting[language=Python]{.././codes-python/line53.py2.py}

The greedy algorithm converges to the theoretical optimum, demonstrating its effectiveness for this problem class. The algorithm's strength lies in its ability to handle discrete constraints and multiple objectives while maintaining computational efficiency.

\subsection{Tier 3: The Advanced Algorithm (Simulated Annealing Metaheuristic)}

For complex real-world scenarios involving multiple products, stochastic demand, or nonlinear cost structures, advanced metaheuristic approaches become essential. Simulated Annealing provides a robust framework for escaping local optima while exploring the solution space efficiently.

\textbf{Simulated Annealing for EOQ with Shortages}

The algorithm mimics the physical annealing process, gradually reducing the ``temperature'' to allow the system to settle into a low-energy state corresponding to the optimal solution.

\textbf{Advanced Algorithm Implementation}

\lstinputlisting[language=Python]{.././codes-python/line53.py3.py}

\textbf{Simulated Annealing Visualization}

\begin{figure}[htbp]
\centering
\includegraphics[width=\textwidth]{../figures-png/line53.fig3.png}
\caption{Simulated Annealing Convergence Pattern}
\end{figure}

\textbf{Advanced Algorithm Results and Analysis}

The Simulated Annealing algorithm provides several advantages over classical approaches:

\textbf{Solution Quality:} The algorithm consistently finds solutions within 0.1\% of the theoretical optimum across multiple runs.

\textbf{Robustness:} Multi-restart capability ensures reliable performance even with suboptimal parameter settings.

\textbf{Flexibility:} The framework easily extends to handle additional constraints such as supplier capacity limits, quantity discounts, or stochastic demand patterns.

\textbf{Example Results for Meridian Electronics:}

\lstinputlisting[language=Python]{.././codes-python/line53.py4.py}

The Simulated Annealing approach demonstrates exceptional consistency, with all five restarts converging to solutions within \$0.12 of the theoretical optimum. This robustness makes it particularly valuable for real-world implementations where solution quality and reliability are paramount.

\section{Practice Questions}

\textbf{Tier 1 Questions (Mathematical Formulation):}

1. \textbf{Multi-Product EOQ with Shortages:} A pharmaceutical distributor manages three critical medications with the following parameters: Product A (D=2400 units/year, K=\$80, h=\$15/unit/year, p=\$25/unit/year), Product B (D=1800 units/year, K=\$60, h=\$12/unit/year, p=\$20/unit/year), and Product C (D=3600 units/year, K=\$100, h=\$18/unit/year, p=\$30/unit/year). Formulate the mathematical model to minimize total system cost across all products, considering a storage capacity constraint of 2000 units total maximum inventory. Calculate the optimal order quantities and shortage levels.

2. \textbf{EOQ with Quantity Discounts and Shortages:} An automotive parts supplier offers a 5\% discount on orders exceeding 2000 units and a 10\% discount on orders exceeding 4000 units. Given D=18000 units/year, K=\$200, h=\$10/unit/year, p=\$15/unit/year, and unit cost=\$50, develop the mathematical formulation to determine the optimal ordering policy that considers both quantity discounts and planned shortages.

\textbf{Tier 2 Questions (Heuristic Implementation):}

3. \textbf{Dynamic Demand EOQ with Shortages:} Implement a greedy heuristic algorithm for a seasonal product with quarterly demand pattern: Q1=8000, Q2=12000, Q3=6000, Q4=10000 units. Order cost K=\$150, holding cost h=\$8/unit/year, shortage cost p=\$12/unit/year. The algorithm should determine optimal ordering schedule and shortage management across the four quarters. Provide the complete Python implementation with example output showing quarterly inventory levels and costs.

4. \textbf{Stochastic EOQ with Service Level Constraints:} Design a heuristic algorithm for an EOQ system where demand follows a normal distribution (mean=3000 units/month, std=500 units/month), K=\$120, h=\$6/unit/year, p=\$18/unit/year, with a minimum service level requirement of 95\%. The algorithm should balance expected shortage costs with service level penalties. Include safety stock calculations and demonstrate convergence behavior.

\textbf{Tier 3 Questions (Advanced Metaheuristic):}

5. \textbf{Multi-Objective EOQ Optimization:} Develop a Simulated Annealing algorithm that simultaneously optimizes three objectives: minimize total cost, minimize maximum shortage level, and maximize service level. Use the weighted sum approach with weights w1=0.6 (cost), w2=0.3 (shortage), w3=0.1 (service). Given D=24000 units/year, K=\$180, h=\$10/unit/year, p=\$20/unit/year, implement the complete algorithm and demonstrate Pareto frontier exploration through different weight combinations.

6. \textbf{Supply Chain Network EOQ with Shortages:} Create an advanced metaheuristic for a three-tier supply chain (supplier-distributor-retailer) where shortages at any level affect downstream operations. Each tier has different cost structures: Supplier (K=\$500, h=\$5/unit/year, p=\$30/unit/year), Distributor (K=\$200, h=\$8/unit/year, p=\$25/unit/year), Retailer (K=\$80, h=\$12/unit/year, p=\$40/unit/year). End customer demand=36000 units/year. Implement a hybrid Simulated Annealing-Genetic Algorithm that optimizes the entire network while considering shortage propagation effects and transportation delays.

\end{document}
